% Source: https://artofproblemsolving.com/community/c3h2604490_geometry_3
% Stable identity: aops-topic-2604490-q08
Always $|AP-PB|\ge0$, with equality exactly when $PA=PB$, i.e. when $P$ lies on the perpendicular bisector $m$ of $AB$.

The line $m$ is perpendicular to $AB$. If $m$ were parallel to $\ell$, then $AB$ would be perpendicular to $\ell$, which is excluded. So $m$ meets $\ell$ in exactly one point $P$, and there $|AP-PB|=0$.

So in both cases (same side or opposite sides) the answer is the same:
\[\boxed{P=\ell\cap m,\ \text{where }m\text{ is the perpendicular bisector of }AB,\ \text{and the minimum is }0.}\]

\begin{asy}
usepackage("amsmath");
size(4.2cm);
defaultpen(fontsize(10pt));
draw((-1.5,0)--(4.5,0)); label("$\ell$",(4.5,0),E);
pair A=(0,2),B=(3,1),M=(A+B)/2,P=(1,0);
draw(A--B,linewidth(0.8)); draw(P-0.25*(1,3)--M+0.5*(1,3),gray(0.4)); label("$m$",M+0.5*(1,3),N,fontsize(9pt));
draw(A--P--B,dashed); dot(A);dot(B);dot(P);dot(M);
label("$A$",A,N);label("$B$",B,E);label("$P$",P,SE);
\end{asy}

\begin{asy}
usepackage("amsmath");
size(4.2cm);
defaultpen(fontsize(10pt));
draw((-1.5,0)--(4.5,0)); label("$\ell$",(4.5,0),E);
pair A=(0,2),B=(3,-1),M=(A+B)/2,P=(1,0);
draw(A--B,linewidth(0.8)); draw(P-0.9*(1,1)--M+1.4*(1,1),gray(0.4)); label("$m$",M+1.4*(1,1),NE,fontsize(9pt));
draw(A--P--B,dashed); dot(A);dot(B);dot(P);dot(M);
label("$A$",A,N);label("$B$",B,S);label("$P$",P,1.6*S);
\end{asy}

Left: same side. Right: opposite sides. In both, $P$ is where the perpendicular bisector $m$ meets $\ell$.
